\section{The Return of Christ and the General Judgment}

No man and no angel knows the date of the Last Day. ``Of that day or hour no man knoweth, neither the angels in heaven, nor the Son, but the Father'' (Mark 13:32; cf.\ Acts 1:7); ``the day of the Lord shall so come as a thief in the night'' (1 Thess 5:2). The \textit{Supplement} observes that because the world ends by divine power alone, knowledge of its hour belongs fittingly to God alone; citing St.\ Augustine's letter to Hesychius (\textit{Epistle} 199), it warns against fixing any month, year, or century from earthly calamities and repeats Augustine's practical rule: ``as each one's last day finds him, so will the world's last day find him.''\endnote{\textit{Supplementum}, q.\ 88, a.\ 3, corpus, obj.\ 4, and ad 2.}

Yet Scripture discloses both the signs preceding the consummation---which the \textit{Roman Catechism} sums up as the preaching of the Gospel throughout the world (Matt 24:14), a great revolt from the faith, and the revelation of Antichrist (2 Thess 2:3)\endnote{\textit{Catechism of the Council of Trent}, Part I, Creed Article VII, trans.\ Donovan, p.\ 64.}---and the majesty of Christ's return:

\begin{quote}
And immediately after the tribulation of those days, the sun shall be darkened and the moon shall not give her light and the stars shall fall from heaven and the powers of heaven shall be moved. And then shall appear the sign of the Son of man in heaven. And then shall all tribes of the earth mourn: and they shall see the Son of man coming in the clouds of heaven with much power and majesty. And he shall send his angels with a trumpet and a great voice: and they shall gather together his elect from the four winds, from the farthest parts of the heavens to the utmost bounds of them. (Matt 24:29--31)
\end{quote}

St.\ Paul adds the archangel's summons: ``For the Lord himself shall come down from heaven with commandment and with the voice of an archangel and with the trumpet of God: and the dead who are in Christ shall rise first. Then we who are alive, who are left, shall be taken up together with them in the clouds to meet Christ, into the air'' (1 Thess 4:15--16). As the Fourth Lateran Council defines, all will rise ``with their own proper bodies which they now bear, that they may receive according to their works, whether they were good or evil''---\textit{cum suis propriis resurgent corporibus, quae nunc gestant}.\endnote{Fourth Lateran Council (1215), Constitution 1, \textit{Firmiter credimus} (DS 801).}

Before the King shines ``the sign of the Son of man,'' unanimously recognized by the Fathers as the Cross. St.\ Cyril of Jerusalem teaches his catechumens:

\begin{quote}
Now Christ's own true sign is the Cross; a sign of a luminous Cross shall go before the King, plainly declaring Him who was formerly crucified \dots\ The sign of the Cross shall be a terror to His foes; but joy to His friends who have believed in Him, or preached Him, or suffered for His sake.\endnote{St.\ Cyril of Jerusalem, \textit{Catechetical Lectures} XV.22, in \textit{NPNF} II, vol.\ 7.}
\end{quote}

The \textit{Supplement} notes that this luminous Cross signifies no present weakness, ``but so as to show how justly those were condemned who scorned so great mercy,''\endnote{\textit{Supplementum}, q.\ 90, a.\ 2, ad 2.} and St.\ Francis de Sales calls it ``a standard of mercy to the good, and of rigor to the wicked.''\endnote{St.\ Francis de Sales, \textit{Introduction to a Devout Life} I.14 (Sixth Meditation: On Judgment), New York: Fr.\ Pustet, p.\ 35.}

Then follows the universal assize foreshown to Daniel---``the judgment sat, and the books were opened'' (Dan 7:10)---and revealed to St.\ John:

\begin{quote}
And I saw a great white throne and one sitting upon it, from whose face the earth and heaven fled away: and there was no place found for them. And I saw the dead, great and small, standing in the presence of the throne. And the books were opened: and another book was opened, which is the book of life. And the dead were judged by those things which were written in the books, according to their works. And the sea gave up the dead that were in it: and death and hell gave up their dead that were in them. And they were judged, every one according to their works. (Apoc 20:11--13)
\end{quote}

In Book XX of \textit{The City of God}, St.\ Augustine dismisses any crudely material parchment scroll and supplies the classical Catholic exposition of those opened books:

\begin{quote}
If this book be materially considered, who can reckon its size or length, or the time it would take to read a book in which the whole life of every man is recorded? \dots\ We must therefore understand it of a certain divine power, by which it shall be brought about that every one shall recall to memory all his own works, whether good or evil, and shall mentally survey them with a marvellous rapidity, so that this knowledge will either accuse or excuse conscience, and thus all and each shall be simultaneously judged.\endnote{St.\ Augustine, \textit{The City of God} XX.14, trans.\ Dods (1871).}
\end{quote}

The Book of Life, Augustine adds, is not an aid to God's memory, but ``His infallible prescience \dots\ in which they are written, that is to say, known beforehand.''\endnote{St.\ Augustine, \textit{The City of God} XX.15, trans.\ Dods (1871).} St.\ Gregory the Great offers a complementary contemplation in his \textit{Moralia in Job}: ``For the Book of Life is the very sight of the approaching Judge. In this are written as it were, all His commands, for whoever beholds it, soon understands by the testimony of conscience what he has omitted to do,'' while the opened books also signify the lives of the just, in whom heaven's precepts stand visibly impressed in act.\endnote{St.\ Gregory the Great, \textit{Moralia in Job} XXIV.8.16, Library of the Fathers translation.}

The \textit{Supplement} develops this illumination with striking balance. Each person will behold both his own deeds and those of all others, so that the equity of every verdict stands manifest; yet the forgiven sins of the saints bring them no shame, ``as neither does it bring confusion to Mary Magdalen that her sins are publicly recalled in the Church,'' turning instead into glory for their penance and praise for God's mercy.\endnote{\textit{Supplementum}, q.\ 87, aa.\ 1--3; a.\ 2, ad 3.} On the outward mode and place of the judgment, the \textit{Supplement} preserves sober restraint: whether the inquiry occurs vocally or mentally ``is not possible to come to any certain conclusion,'' though ``it is, however, the more probable opinion that the whole of this judgment \dots\ will take place mentally''; likewise, regarding the ancient tradition drawn from Joel 3:2 and Acts 1:11 placing Christ's descent near the Valley of Josaphat and Mount Olivet, ``we cannot know with any great certainty the manner in which this judgment will take place.''\endnote{\textit{Supplementum}, q.\ 88, aa.\ 2, 4.}

Around Christ's throne sit His saints as judicial assessors. Our Lord promised the Apostles, ``you also shall sit on twelve seats judging the twelve tribes of Israel'' (Matt 19:28), and St.\ Paul asks, ``Know you not that the saints shall judge this world? \dots\ Know you not that we shall judge angels?'' (1 Cor 6:2--3). St.\ Augustine explains that twelve symbolizes the entire company of the saints who judge with Christ---else St.\ Paul, who labored more than all, would lack a seat---and the twelve tribes signify all nations.\endnote{St.\ Augustine, \textit{The City of God} XX.5, trans.\ Dods (1871).} The \textit{Supplement} clarifies that authoritative dominion belongs to Christ alone; the saints, especially the Apostles and those who left all to follow Him, judge assessorially by approving His sentence and manifesting divine justice, while the angels serve as witnesses and ministers.\endnote{\textit{Supplementum}, q.\ 89, aa.\ 1--3.}

Before that throne every earthly rank dissolves. St.\ Cyril of Jerusalem comforts the poor and forgotten:

\begin{quote}
Be of good courage, O man; the Judge is no respecter of persons \dots\ He honours not the learned before the simple, nor the rich before the needy. \dots\ Though thou be a slave, though thou be poor, be not any whit distressed; He who took the form of a servant despises not servants. \dots\ Every prayer of thine, every Psalm thou singest is recorded; every alms-deed, every fast is recorded.\endnote{St.\ Cyril of Jerusalem, \textit{Catechetical Lectures} XV.23, in \textit{NPNF} II, vol.\ 7.}
\end{quote}

St.\ Alphonsus Liguori contrasts the despised saints and martyrs shining in Christ's court with Herod, Pilate, and Nero stripped of earthly pomp, ``who, for the sake of making a figure for a short time on the theatre of this world, shall afterwards have to act the part of reprobates in the tragedy of judgment.''\endnote{St.\ Alphonsus Maria de Liguori, \textit{Preparation for Death}, Consideration XXV, Second Point (Dublin: James Duffy, 1843), p.\ 306.}

At the center of the General Judgment stands Matthew 25:31--46:

\begin{quote}
And when the Son of man shall come in his majesty, and all the angels with him, then shall he sit upon the seat of his majesty. And all nations shall be gathered together before him: and he shall separate them one from another, as the shepherd separateth the sheep from the goats: And he shall set the sheep on his right hand, but the goats on his left. Then shall the king say to them that shall be on his right hand: Come, ye blessed of my Father, possess you the kingdom prepared for you from the foundation of the world. For I was hungry, and you gave me to eat: I was thirsty, and you gave me to drink: I was a stranger, and you took me in: Naked, and you covered me: sick, and you visited me: I was in prison, and you came to me. \dots\ Amen I say to you, as long as you did it to one of these my least brethren, you did it to me. (Matt 25:31--36, 40)
\end{quote}

By the same concrete measure of charity withheld from the hungry, the stranger, the sick, and the imprisoned, the King passes sentence on His left hand, closing history with the eternal parting: ``And these shall go into everlasting punishment: but the just, into life everlasting'' (Matt 25:46). St.\ Cyril of Jerusalem cuts off every evasion: the shepherd distinguishes sheep from goats by evident marks, and ``there is here no need of allegory, but to fulfil what is said.''\endnote{St.\ Cyril of Jerusalem, \textit{Catechetical Lectures} XV.25--26, in \textit{NPNF} II, vol.\ 7.} In the teaching of the \textit{Catechism of the Catholic Church}, before Christ who is Truth itself each person's true standing with God will be laid bare, unveiling to its furthest consequences all the good each has done or failed to do during earthly life.\endnote{\textit{Catechism of the Catholic Church}, \S\S 1038--1039.}

After that sentence no intermediate state remains: Purgatory is ended, death and hell have surrendered their dead, and only the eternal kingdom and the eternal abyss abide.
